\section{Zielschema}\label{sec:zielschema}

% Zielumfang: ca. 1 Seite. Schmales Sternschema; Entscheidungen begründen, Alternativen nennen.

\todo{Abbildung: Sternschema DIM\_KUNDE, DIM\_KUNDE\_ADRESSE (historisiert), DIM\_VERKAEUFER, DIM\_PRODUKT, FAKT\_VERKAUF, MAP\_SCHLUESSEL. Als Screenshot aus Excel/Draw.io oder TikZ.}

\begin{figure}[H]
  \centering
  % \includegraphics[width=0.9\textwidth]{zielschema}
  \fbox{\parbox{0.9\textwidth}{\centering\vspace{2cm}\todo{Abbildung Zielschema}\vspace{2cm}}}
  \caption{Zielschema des integrierten Datenbestands}
  \label{fig:zielschema}
\end{figure}

\paragraph{DIM\_KUNDE.} Surrogatschlüssel \pq{KundeKey} als Index \skript{8.2.2.3}, daneben die drei Quellschlüssel \pq{ERP\_KdNr}, \pq{Shop\_ID}, \pq{CRM\_ID} (leer, wo nicht vorhanden), \pq{Kundentyp} auf vier Zielwerte, \pq{Land} als ISO-2, \pq{Newsletter} als Boolean, \pq{Herkunftsquellen} zur Nachvollziehbarkeit.

\paragraph{DIM\_KUNDE\_ADRESSE.} Adressen werden bei Widerspruch zwischen ERP und CRM nicht entschieden, sondern historisiert \skript{8.5}: \pq{GueltigVon}, \pq{GueltigBis}, \pq{CurrentFlag}. Ein Auftrag wird der zum Auftragsdatum gültigen Adresse und damit Region zugeordnet. \todo{Begründung: verhindert die falsche Verlagerung des Einkaufsverhaltens bei Umzug. Alternative \enquote{jüngste gewinnt} nennen und verwerfen.}

\paragraph{DIM\_VERKAEUFER.} \pq{VertreterNr}, normalisierter Name, Rolle (Außendienst, Innendienst, \enquote{Online} als Pseudo-Verkäufer für Shop-Direktkäufe), Region, Aktiv-Kennzeichen. Der ausgeschiedene Vertreter V01 bleibt mit Austrittsdatum erhalten, weil seine Aufträge sonst ihre Referenz verlieren.

\paragraph{DIM\_PRODUKT.} \pq{ERP\_ArtNr}, \pq{Shop\_SKU}, Produktgruppe -- aus dem sprechenden ERP-Schlüssel bzw. aus \quelle{products.category} abgeleitet.

\paragraph{FAKT\_VERKAUF.} Eine Zeile je Auftragsposition: Datum (Jahr, Quartal, Saison), die drei Fremdschlüssel, \pq{Kanal} (Direkt/Online), Menge in Stück, \pq{UmsatzNettoEUR} (CHF umgerechnet, Brutto je Land herausgerechnet), \pq{Lieferdauer\_Tage}, \pq{Zahlungsverzug\_Tage}, \pq{Quelle}. Shop-Aufträge sind auch im ERP importiert; Umsatz wird deshalb nur aus einer Quelle übernommen. \todo{Entscheidung begründen, welche.}

\paragraph{Entscheidungen.} Leads kommen als Kunden ohne Umsatz mit Kontaktstatus \enquote{Lead} ins Ziel (Vollständigkeit; Alternative: filtern). Die Lieferart wird Attribut, nicht Entität (drei Werte ohne eigene Eigenschaften, Abb.~8.10 im Skript). Der Verkäufer hängt am Auftrag, nicht am Kunden; der Stammdaten-Vertreter ist nur Default für Aufträge ohne Angabe -- nur so ist Frage F5 (Kanal) korrekt beantwortbar \skript{8.2.2.5}.

\Cref{tab:fragencheck} prüft, dass alle sechs Auswertungsfragen aus dem Schema beantwortbar sind.

\begin{table}[H]
  \centering\small
  \caption{Abdeckung der Auswertungsfragen durch das Zielschema}
  \label{tab:fragencheck}
  \begin{tabularx}{\textwidth}{@{}>{\raggedright\arraybackslash}p{4.2cm} >{\raggedright\arraybackslash}p{4.6cm} X@{}}
    \toprule
    Frage & Fakt $\times$ Dimension & Benötigte Integrationsleistung \\
    \midrule
    F1 Umsatz je Verkäufer $\times$ Kundentyp & FAKT $\times$ VERKAEUFER $\times$ KUNDE & Kundentyp-Mapping (3 Kodierungen) \\
    F2 $\times$ Land/Region & $\times$ KUNDE\_ADRESSE & Länder-Mapping, Region aus PLZ, Historisierung \\
    F3 $\times$ Saison & Datumsattribute & Datumsformate, Zeitzone \\
    F4 $\times$ Produktgruppe & $\times$ PRODUKT & ArtNr$\leftrightarrow$SKU-Mapping \\
    F5 Kanalanteil & Kanal, Verkäufer $\neq$ Online & Record Linkage Shop$\leftrightarrow$ERP/CRM \\
    F6 Zahlungs-/Lieferqualität & Lieferdauer, Zahlungsverzug & Zahlungsziel-Defaults, Nulldatum \\
    \bottomrule
  \end{tabularx}
\end{table}
